\chapter{Inspecting, Debugging, Filing}
\label{ch:tools}

\standfirst{The Browser has a chapter of its own. These are the other four
tools you live with --- all written in the language, all changeable from
inside themselves.}

\section{The Workspace}

A text pane with no model behind it. Type expressions, select them, \ct{do it}
or \ct{print it}. Variables you assign to without declaring become workspace
variables and persist while the window is open, which makes it a scratch pad
rather than a series of one-shot evaluations.

\ct{system workspace} on the desktop menu opens one pre-filled with useful
expressions to select and run.

\section{The Inspector}

Send \ct{inspect} to anything:

\begin{st}
(1 to: 10) asOrderedCollection inspect
\end{st}

A two-pane window: instance variables on the left, the value of the selected
one on the right. \ct{self} is the first entry. The right-hand pane is a text
editor with \ct{do it} and \ct{print it}, and inside it \ct{self} is the
object you are inspecting---so you can poke at a live object with real code:

\begin{st}
self addFirst: 0
self size
\end{st}

The middle button in the left pane offers \ct{inspect} again, which opens an
Inspector on
the selected field. That is how you walk into a structure.

\section{The Debugger}

When something goes wrong and nothing handles it, a \emph{notifier} appears: a
small window with the error and the top few frames.

\begin{center}
\small
\begin{tabular}{@{}lp{0.60\textwidth}@{}}
\toprule
\ctp{proceed} & carry on as if nothing happened\\
\ctp{abandon} & give up on this process\\
\ctp{debug} & open the full Debugger\\
\bottomrule
\end{tabular}
\end{center}

\ct{debug} gives you a window with the call stack at the top, the source of
the selected frame in the middle, and two Inspectors at the bottom---one on
the receiver, one on the method's own variables.

What makes this different from a debugger in another language:

\begin{itemize}
\item The source pane is \emph{editable}. Fix the method, \ct{accept}, and the
  frame is restarted with the new code.
\item The receiver Inspector is live. Set the instance variable that was
  \ct{nil} and carry on.
\item \ct{proceed} resumes from where the error was raised, so once you have
  fixed the cause you continue rather than starting over.
\end{itemize}

This is the loop the language was designed around: hit an error, fix it in
place, continue. It repays getting comfortable with.

\begin{stnote}[The debugger needs a screen]
A headless run---\ct{-eval}, \ct{-tests}, \ct{-doctests}---reports errors as
text on standard error instead, because there is nowhere to draw. That is
deliberate; an error during a bootstrap must not try to open a window.
\end{stnote}

\section{The File List}
\label{sec:filelist}

\ct{file list} on the desktop menu. The top pane takes a pattern; type one and
\ct{accept} it. The pattern is a shell-style glob:

\begin{st}
*
*.st
/home/blake/*.txt
\end{st}

The list fills with matching names. Selecting one and choosing
\ct|get contents| from the middle-button menu shows it in the bottom pane. The
file menu is \ct{get contents}, \ct{file in}, \ct{copy name}, \ct{rename},
\ct{remove}.

\ct{file in} reads a chunk-format source file and installs everything in
it---which is how code moved between images before there were packages, and
still works.

\subsection{Directories}

\ct|get contents| on a \emph{directory} enters it rather than trying to read
it: the list becomes that directory's contents and the pattern pane says where
you now are. 1983 had no word for a directory at all and simply failed here.

The pattern is both the display and the control, so moving around is always
the same gesture --- edit the top pane and \ct{accept}:

\begin{center}
\small
\begin{tabular}{@{}lp{0.58\textwidth}@{}}
\toprule
\ctp{doc/*} & what entering \ctp{doc} writes for you\\
\ctp{*} & back to where the system was started\\
\ctp{../*} & above it\\
\bottomrule
\end{tabular}
\end{center}

Names in the list carry their directory with them, so what you select is what
you get however deep you have gone.

\section{Finding things without the Browser}
\label{sec:withoutbrowser}

\begin{st}
Smalltalk at: #OrderedCollection      "the class, by name"
Smalltalk includesKey: #Foo
Object allSubclasses size
Collection allSubclasses
OrderedCollection selectors
OrderedCollection allInstVarNames
OrderedCollection comment
(OrderedCollection >> #add:) getSource
\end{st}

And from the command line, without opening a window at all:

\begin{sh}
./st80 -classes st80.image | grep -i collection
./st80 -methods st80.image | grep 'inject:into:'
./st80 -disasm st80.image OrderedCollection add:
\end{sh}
